\section{Everyone Else Knows Better}

Deriving some insights from \emph{The Crisis of the Modern World}, by \textbf{Rene Guenon}. 

\paragraph{Rationalism}
Rationalism is the notion that human reason alone — by its ability to categorize, analyze, discuss, compare, contrast, interpret, etc. — is an adequate instrument to determine the truth. Guenon points out that is is a modern phenomenon.

\begin{quotex}
rationalism is a specifically modern phenomenon, one that is closely connected with individualism, being nothing other than the negation of any faculty of a supra-individual order.

\end{quotex}
The Enlightenment liberal believes that if the disputing parties engage in a rational discussion of ideas, then they will come to a satisfactory conclusion. \textbf{However, everyone else knows better}.

It is obvious, by now, that there is no rational answer to the ultimate questions of life; discussions and debates will continue interminably with no resolution. There is the obvious problem associated with individualism and its corollary egalitarianism. In point of fact, not everyone is competent to even understand, never mind debate the issues. Not even science can help.

The conclusions are few:

\begin{itemize}
\item Accept that there is no generally agreed upon solution, so that it can be no more than a personal choice. And like the Highlanders, the logical battle of ideas continues on indefinitely. 
\item There must be a higher form of knowledge that transcends the limits of human reason, and enters into the human world from outside the system. 
\end{itemize}
\paragraph{Intuition}
The intuitive grasp of principles, which is a direct understanding not mediated through the methods of rationalism, is the solution to the artificial problems posed by rationalism. The difference is that intuition cannot be simply taught since it requires the union of the knower and the known, which is experiential. Unlike rational, which requires nothing of its adherents, intuition requires a change in one's level of being. Traditional truths are immutable and once understood, they are seen to be infallible.

\begin{quotex}
Intellectual intuition, by which alone metaphysical knowledge is to be obtained, has absolutely nothing in common with this other `intuition' of which certain contemporary philosophers speak: the latter pertains to the sensible realm and in fact is sub-rational, whereas the former, which is pure intelligence, is on the contrary supra-rational. But the moderns, knowing nothing higher than reason in the order of intelligence, do not even conceive of the possibility of intellectual intuition, whereas the doctrines of the ancient world and of the Middle Ages explicitly recognized its existence and its supremacy over all the other faculties.

Intuition is immutable and infallible in itself, and the only starting-point for any development in conformity with traditional norms.

\end{quotex}
\paragraph{Verbal Disputes}
Rationalism tries to sneak into the discussion. As Guenon emphasizes in this passage, metaphysics does not engage in philosophical arguments. Trying to capture ultimate truths verbally is a fool's errand. A fortiori, verbal disputes about the meanings of words are a complete waste of time. Knowledge of principles does not change by verbal formulations.

\begin{quotex}
It should be clearly understood that it would be utterly useless to put forward here, by way of objection, any more or less specious philosophical arguments; we are speaking seriously, of serious matters, and have no time to waste on verbal disputes that would be of no interest, and could serve no useful purpose.

\end{quotex}
\paragraph{The Party Line}
Guenon makes perfectly clear that he is not interested in any sort of philosophical ``school" nor in party politics. Moreover, Western labels like pantheism or idealism are utterly inapplicable to metaphysical teaching.

\begin{quotex}
It is our intention to remain entirely aloof from all controversies and quarrels of \textbf{school or party}, just as we refuse absolutely to accept any Western label or definition, since none is applicable; whether this is pleasing or displeasing, it is a fact, and nothing will change our attitude in this regard.

\end{quotex}
\paragraph{Warning to the Possible Elite}
``Many are called but few are chosen." There are many who have a vague interest in Tradition; they are the called. But the chosen, who have achieved a higher realization, are few. The called only have the possibility of knowledge. The chosen are not seduced by modern ideas, but the called are vulnerable. They need to be extra vigilant is they truly desire to attain any degree of realization.

\begin{quotex}
A warning must be addressed to those who, because of their capacity for a higher understanding, if not because of the degree of knowledge to which they have actually attained, seem destined to become elements of a possible elite. There is no doubt that the force of modernism, which is truly `diabolic' in every sense of the word, strives by every means in its power to prevent these elements, today isolated and scattered, from achieving the cohesion that is necessary if they are to exert any real influence on the general mentality. It is therefore for those who have already more or less completely become aware of the end toward which their efforts should be directed to stand firm against whatever difficulties may arise in their path and threaten to turn them aside.

\end{quotex}


\flrightit{Posted on 2022-01-18 by Cologero }

\begin{center}* * *\end{center}

\begin{footnotesize}\begin{sffamily}



\texttt{Mikail on 2022-01-19 at 06:53 said: }

``rationalism is a specifically modern phenomenon, one that is closely connected with individualism, being nothing other than the negation of any faculty of a supra-individual order."

It seems that Protestant thinking is the very instigator of the modern order. One could say that the reformation and the scientific revolution occuring in the same time-frame is purely coincidental, but it seems there is a causal link. So which caused which? Did the scientific revolution cause the reformation, or the other way around? Considering the emphasis on rational thought, individualism and free inquiry among Protestant thinkers, the latter could equally be the case.


\hfill

\texttt{Mike M on 2022-01-19 at 11:44 said: }

``It seems that Protestant thinking is the very instigator of the modern order…be the case."

On the level of earthly manifestation that is a certain way to investigate what occurred. However, causes occur in reality before they manifest, ie. they don't come from nowhere. If Protestant thinking is the instigator, who instigated the instigator? And so on, if infinite regress is an issue, that typically points to a cause outside of the rational sphere. It may be helpful to meditate on why both came about and what the goal is. Barfield's Saving the Appearances may help you if you are truly interested in this question.


\hfill

\texttt{Santiago on 2022-01-19 at 18:23 said: }

Great post. According to Barfield, the idolatry is still a necessary step. The iguana can't find food without scuffling around like a retarded inbred first. Either it gives up, or it doesn't.

As for the labels, ``idealism", ``moderate realism", ``monophysitism", ``Ayn Randian Objectivism", etc. I can see how it's easy to fall into these ideological camps which can taint an entire worldview, as they place an undue emphasis on systems, rather than what is being sought.


\hfill

\texttt{Cassiodorus on 2022-01-29 at 09:15 said: }

The notion that `` Tradition" isn't philosophy is highly dubious . The doctrine of unqualified non- dualism, maya, the degrees of reality and the like are all philosophical doctrines 

Gnosis in a traditional Catholic context is light years away from the gnosis of Advaita/ Perrenialism. The spiritual experiences might actually be the same , but their interpretations vary.

What could possibly make one interpretation higher than another if not philosophy ? But that would contradict Guenon's main point about `` profane philosophy".


\hfill

\texttt{Cologero on 2022-01-29 at 11:56 said: }

How many years, Cassiodorus, have you been reading and commenting here?

Well, you finally nailed it.

As you say, the philosopher is interested in ``interpreting" spiritual experiences which he himself never experienced.

There are only a very few who pursue the ``experience", or better said, intellectual intuition.


\hfill


\end{sffamily}\end{footnotesize}
